\section{Princípios da estimação do \textit{Sound Speed Field}}

Embora a velocidade do som na água seja considerada constante em torno de $1500 \si{\meter}/\si{\second}$, ela pode variar em ambientes de pequena e grande escala, respectivamente, em proporções pequenas e grandes \cite{kieffer_sound_1977,trusler_determination_2017,li_sound_2015,pinson_sound_2010}. Muitas propriedades podem afetar a velocidade do som na água, como pressão, temperatura, salinidade e correntes marítimas \cite{wilson_speed_1959,medwin_fundamentals_1999,mcdougall_getting_2011,delgrosso_new_1974}. À medida que essas propriedades variam na massa de água, a velocidade também varia, criando um campo de velocidade do som (SSF, do inglês \textit{Sound Speed Field}) na massa de água. Embora o termo perfil de velocidade do som (SSP) \cite{li_acoustic_2019,li_sound_2015,pinson_sound_2010} também seja comumente usado, ele geralmente se refere ao cenário de estimativa 1D.

Os trabalhos iniciais sobre a estimativa do campo de velocidade na água visavam estimar essas propriedades principais, bem como a relação entre elas e a velocidade da água. Trabalhos como \cite{wilson_speed_1959,medwin_fundamentals_1999,mcdougall_getting_2011,delgrosso_new_1974,chen_speed_1977,wong_speed_1995,roquet_accurate_2015} operaram dentro dessa estrutura, e alguns ainda são usados hoje para modelar o comportamento de correntes de água em escala global.

Em \cite{wilson_speed_1959}, experimentos mostraram que (em água destilada) $c$ tem uma relação estritamente positiva com a pressão e uma correlação parabólica com a temperatura. Ou seja, um aumento na pressão corresponde diretamente a um aumento em $c$, e similarmente com a temperatura até um certo ponto, após o qual aumentos na temperatura ou não afetam $c$, ou até mesmo causam uma diminuição na velocidade. Em \cite{medwin_fundamentals_1999}, são apresentados modelos empíricos para a velocidade do som (c), que também consideram a salinidade em sua formulação.

O modelo TEOS-10 \cite{mcdougall_getting_2011} é outro método para estimar o campo de velocidade, dados os parâmetros de temperatura, pressão e salinidade na massa de água, resolvendo equações químico-termodinâmicas para encontrar a distribuição de velocidade adequada.

A busca pela velocidade do som em cada ponto de uma massa de água --- a velocidade do som em cada ponto --- permite o conhecimento das propriedades de uma onda que se propaga nesse meio. Pela lei de Snell, uma mudança na velocidade de uma onda também muda sua direção; portanto, uma onda que se propaga em uma massa de água suficientemente grande também será afetada dentro dela, resultando em curvas, refrações, reflexões e formação/dissolução de feixes.

Conhecendo a velocidade do som em cada ponto, é possível saber a trajetória que uma onda unidirecional percorreria. Ao tratar uma onda omnidirecional como um conjunto de ondas unidirecionais, é possível calcular o tempo de propagação entre a fonte e qualquer ponto dentro do meio, bem como o perfil da sua frente de onda.

A maioria das abordagens modernas para estimativa do SSF evita estimar essas propriedades, e trata essa estimatição como um problema de inversão. Ou seja, como o conhecimento da SSF permite calcular o caminho de uma onda através do meio, o conhecimento desses caminhos permite estimar a SSF no meio. Diferentes técnicas tratam esse problema de inversão usando diferentes tipos de dados e abordagens matemáticas diferentes.

A maioria das técnicas utiliza fontes e receptores sonoros dentro da massa de água e, comparando as propriedades das ondas produzidas e recebidas, estima o caminho das ondas, disso estimando a SSF da massa de água.

\subsection{Métodos para inversão do SSF}

Diversos métodos podem ser utilizados para estimar o SSF, usando diferentes informações adquiridas por sensores ao longo da massa de água. Esses dados podem incluir o tempo de trânsito entre a fonte e os receptores, a pressão sonora medida, a diferença de fase e amplitude entre os dispositivos, ou mesmo a temperatura, a salinidade e a pressão utilizadas nos métodos citados anteriormente. Além disso, as técnicas podem utilizar essas medições por meio de métodos diretos de estimativa da velocidade do som no mar, bem como por meio de aprendizado de máquina e técnicas estatísticas.

\subsubsection{\textit{Travel-Time Tomography}}

A \textit{Travel-Time Tomography} (TTT) \cite{bormann_seismic_2012,stefanov_travel_2019} é um dos métodos mais simples para estimativa do SSF. Ela utiliza a diferença no tempo de chegada dos sinais acústicos nos sensores, exigindo uma distribuição espacial de fontes e receptores para resolução espacial suficiente.

Em geral, o SSF é discretizado em camadas ou células, onde cada uma é tratada tendo uma velocidade do som constante. Em seguida, calcula-se o comprimento do percurso da onda em cada camada, resultando em um sistema de equações usado para estimar a velocidade em cada célula. Também pode funcionar com funções assumidas para a SSF, descartando a suposição de camadas, mas aumentando a complexidade do problema.

Sua simplicidade, robustez e eficiência computacional são suas principais vantagens, sendo esta última capaz de permitir a implementação em tempo real. No entanto, depende de um posicionamento e sincronização precisos dos dispositivos, e sua suposição de discretização resulta em resolução limitada.

\subsubsection{\textit{Matched Field Processing} e {Matched Field Inversion}}
Nas técnicas de \textit{Matched Field Processing} (MFP) \cite{gemba_adaptive_2017} e \textit{Matched Field Inversion} (MFI) \cite{dosso_bayesian_2011}, os dados de pressão sonora medidos nos receptores são comparados a modelos de pressão acústica para diferentes SSFs, adaptando os modelos aos dados por meio de algoritmos de otimização, a fim de encontrar o mais similar.

Essas técnicas são semelhantes à formação de feixe (beamforming), utilizada em diversas áreas como acústica e processamento de sinais \cite{schmidt_environmentally_1990,brienzo_broadband_1993}, onde ambas exploram a correlação espacial entre os receptores e suas medições, modelando o campo acústico através deles.

O MFP e o MFI também empregam a decomposição modal dos sinais, com cada modo se comportando de maneira diferente na água. Essas duas técnicas incorporam informações mais complexas sobre os sinais, como amplitude e fase, dependendo não apenas do tempo de propagação. Entre suas vantagens, destacam-se o melhor aproveitamento dos dados obtidos, a resolução espacial aprimorada (em comparação com o TTT) e a maior robustez ao ruído devido ao processamento coerente. No entanto, requerem informações espaço-temporais sobre a fonte e os receptores, como diretividade e comportamento espectral. Além disso, não são muito adequadas para corpos muito grandes, visto que um corpo maior resultaria em um modelo de pressão acústica exponencialmente mais complexo; e são computacionalmente mais complexas, além de suscetíveis a mínimos locais.

\subsubsection{\textit{Full Waveform Inversion}}

As técnicas de \textit{Full Waveform Inversion} (FWI) \cite{virieux_overview_2009,fichtner_multiscale_2013} visam a reconstrução da forma de onda dos sinais presentes no SSF, considerando diversos comportamentos que são negligenciados pelas técnicas mencionadas anteriormente. Nesse método, o erro entre as formas de onda estimadas e medidas é minimizado, sendo implementado nos domínios do tempo ou da frequência.

Dado isso, a técnica FWI está entre as que melhor utilizam os dados disponíveis dos receptores, resultando no SSF mais fiel, além de aproveitar ao máximo os sinais disponíveis. Entretanto, devido à consideração de interações complexas entre os sinais e múltiplos caminhos de onda, há um aumento significativo na complexidade matemática e computacional. Isso, aliado à não linearidade do problema, resulta em uma técnica muito complexa e, por vezes, inviável, que também requer uma estimativa inicial adequada dos parâmetros para sua robustez e convergência.